\section{Previous Related Work}
{\bf Dataset Distillation.} The work of \cite{Wangetal18} introduced the problem of \emph{distilling} a dataset into a representative (and smaller) synthetic dataset, in the setting of supervised learning. This and other works e.g. \cite{deng2022remember} use a bi-level optimization formulation in which the model is optimized on the training dataset while the synthetic dataset is optimized on the trained model. A related set of methods rely on matching various properties of the synthetic dataset with the training dataset. In particular, the work of \cite{zhao2021dataset,Zhao2021DatasetCW} matches the model's loss gradients on the training and synthetic datasets as an optimization over the latter, while the model is alternately optimized over the training dataset. In a similar vein, the work of \cite{Wangetal22} aligns the feature distribution of the two datasets in a dynamic bi-level optimization approach, while the works of \cite{Cazenavetteetal22,Cuietal23} match the training trajectory of an initial model optimized on the training dataset with its training trajectory on the synthetic dataset, by optimizing on the latter. Unlike the mentioned works optimizing a model along with the synthetic dataset, the work of \cite{ZBDistribMatch} instead matched the feature-distributions on the training and synthetic datasets with respect to a fixed collection of randomly sampled networks. For linear ridge regression, the work of \cite{nguyen2021dataset} implicitly matched the regression losses by minimizing a surrogate objective, while proving convergence of synthetic feature-vectors given the synthetic labels. More recently, \cite{chen2024provable} analytically gave an efficient synthetic dataset generation algorithm for linear ridge regression, requiring however access to the optimal regression model for the training dataset.


Closely related to DD for linear regression is \emph{matrix sketching} which  provides a principled way to reduce the dimensionality (or size) of training data. By applying randomized projections (e.g., the Johnson--Lindenstrauss transform~\citep{johnson1984extensions,sarlos2006improved}) or leverage-score sampling~\citep{drineas2006fast,mahoney2011randomized}, one can construct a much smaller \emph{sketch} of the original data while provably preserving the solution quality of least-squares regression. 

However, these guarantees are largely limited to linear and convex models. In contrast, DD methods can be applied to neural networks, albeit without comparable theoretical guarantees. Additionally, DD techniques optimize the synthetic dataset, and have been shown to to work well in practice on real data in terms of size compression and quality preservation. 


{\bf Offline RL.} 
A key advantage is that 
offline RL is adaptable to data generated by sub-optimal policies~\citep{levine2020offlinereinforcementlearningtutorial,fu2021d4rldatasetsdeepdatadriven}, while also being scalable to large datasets~\citep{Lange2012}. On the other hand, the lack of online exploration can lead to less generalizable policies being trained due to distributional shifts. To address this, several techniques based on constraining the policy \citep{fujimoto2021a,tarasov2023revisiting} or regularization  of $Q$-value predictors have been developed~\citep{Kumaretal20,kostrikov2022offline}. A major practical consideration is the exponential growth in dataset sizes, along with which the associated challenges of storage, transfer and training on datasets as well as maintaining privacy controls have only multiplied. Dataset distillation has been proposed to tackle these problems, with existing works developing behavioral cloning methods for distilling an offline training dataset comprising trajectories of (state, action)-pairs. While \cite{lupu2024behaviour} propose behavior distillation for online RL, the work of \cite{lei2024offline} tackles the offline RL case by first training an action-value predictor and using action-value weighted BC to optimize the synthetic dataset. On the other hand, the method proposed by \cite{light2024datasetdistillationofflinereinforcement} extends the matching loss gradients approach to offline RL. In particular, it optimizes a synthetic dataset to match the behavioral cloning loss gradients on the offline and synthetic data, where the gradients are with respect to the parameters of sampled action predictor networks. 

